\chapter{磁与磁场}\label{ch:g11:magnetic-fields}

千余年来，水手们借助被磁化的指针导航，
却无人知晓指针为何指向北方。1820 年，
一股电流使指南针偏转：磁性来自运动的电荷。
本章描绘磁场——磁铁、地球与电流的磁场——
然后闭合这一回路：磁场对电流施力，
这正是所有电动机转动的原理。

\section{磁铁与磁场}

\begin{definition}[磁铁与磁极]\label{def:g11:magnetic-fields:magnet}
\emph{磁铁}\index{磁铁}能吸引铁，用细线悬挂后，
会转向北方。其作用集中于两个
\emph{磁极}\index{磁极}，即\emph{北极}（指向地理北方的
一端）与\emph{南极}：同名磁极相斥，异名磁极相吸。
把磁铁锯成两半，得到两块完整的磁铁——
至今尚无实验能分离出单独的磁极。
\end{definition}

\begin{definition}[磁场]\label{def:g11:magnetic-fields:bfield}
磁铁改变其周围的空间
（见\cref{ch:g11:electric-gravitational-fields}），在其周围各点产生一个
\emph{磁场}\index{磁场} $\vect B$。指南针可指示其方向
（指针与 $\vect B$ 同向，北端朝前）；
\emph{霍尔探头}\index{霍尔探头}可测量其大小，
单位为\emph{特斯拉}\index{特斯拉}（\unit{T}）。
\end{definition}

\begin{definition}[磁感线]\label{def:g11:magnetic-fields:lines}
\emph{磁感线}\index{磁感线}是一些曲线，
其切线方向与 $\vect B$ 一致，磁场强处线条密集。在磁铁外部，
它们由北极指向南极；与电场线不同，
磁感线没有端点：每一条都在磁铁内部闭合。
\end{definition}

\begin{omfigure}
\begin{tikzpicture}[scale=0.9,
  fline/.style={omThm, postaction={decorate},
    decoration={markings, mark=at position 0.5 with {\arrow{Stealth}}}}]
  \fill[omProp!15] (-1.4,-0.35) rectangle (0,0.35);
  \fill[omDef!15] (0,-0.35) rectangle (1.4,0.35);
  \draw[thick] (-1.4,-0.35) rectangle (1.4,0.35);
  \draw (0,-0.35) -- (0,0.35);
  \node[font=\small] at (-0.7,0) {S};
  \node[font=\small] at (0.7,0) {N};
  \draw[fline] (1.4,0.2) .. controls (2.6,0.9) and (-2.6,0.9) .. (-1.4,0.2);
  \draw[fline] (1.4,0.3) .. controls (3.8,2.0) and (-3.8,2.0) .. (-1.4,0.3);
  \draw[fline] (1.4,-0.2) .. controls (2.6,-0.9) and (-2.6,-0.9) .. (-1.4,-0.2);
  \draw[fline] (1.4,-0.3) .. controls (3.8,-2.0) and (-3.8,-2.0) .. (-1.4,-0.3);
  \draw[fline] (1.4,0) -- (2.95,0);
  \draw[omExo, thick] (0,1.1) circle (0.22);
  \draw[-Stealth, omDef, thick] (0.22,1.1) -- (-0.22,1.1);
  \draw[omExo, thick] (2.45,0) circle (0.22);
  \draw[-Stealth, omDef, thick] (2.23,0) -- (2.67,0);
\end{tikzpicture}

{\small 条形磁铁的磁感线：由 N 极射出，进入 S 极，
经磁铁内部闭合。指南针放在任意位置都沿过该点的磁感线方向排列。}
\end{omfigure}

\section{地球是一块磁铁}

\begin{definition}[地球磁场]\label{def:g11:magnetic-fields:earth}
地球熔融的铁质地核中的电流使地球成为一块巨大的
磁铁：地表处 $B \approx \qty{50}{\micro T}$，可自由转动的
指针沿磁场的水平分量排列，这便是指南针的原理。
地磁极与地理极不重合；指南针北与真北之间的夹角
称为
\emph{磁偏角}\index{磁偏角}，每次定向都需加以修正。
\end{definition}

\begin{example}[方位修正]\label{ex:g11:magnetic-fields:bearing}
当磁偏角为西偏 $\ang{7}$ 时，船只沿“指南针北”
方向航行 \qty{5.0}{km}，将向西偏离航线 $5000 \times \sin\ang{7} \approx
\qty{6.1e2}{m}$。若要沿真北航行，应使航向比指针方向向 $\ang{7}$
\emph{东}偏转。
\end{example}

\section{电流的磁场}

\begin{remark}[奥斯特，1820]\label{rem:g11:magnetic-fields:oersted}
1820 年讲课时，汉斯·克里斯蒂安·奥斯特看到指南针的指针
在他附近闭合电路时发生摆动：电与磁原本属于同一
门学问——\emph{电流产生磁场}。
\end{remark}

\begin{proposition}[直导线的磁场]\label{prop:g11:magnetic-fields:wire}
通有电流 $I$ 的长直导线产生的磁场，
其磁感线是位于垂直导线平面内、以导线为圆心的圆。
距导线 $d$ 处，
\[
B = \frac{\mu_0\, I}{2\pi d},
\qquad \mu_0 = 4\pi\times 10^{-7}\ \unit{T.m/A} \approx \qty{1.26e-6}{T.m/A}.
\]
\end{proposition}
\admitted

\begin{omfigure}
\begin{tikzpicture}[scale=0.85]
  \node[font=\large] at (0,0) {$\odot$};
  \foreach \r in {0.55,1.05,1.55}
    \draw[omThm,-Stealth] (\r,0) arc (0:350:\r);
  \node[omThm, font=\small] at (45:1.9) {$\vect B$};
  \node[font=\small] at (0,-2.1) {电流朝向读者};
\end{tikzpicture}\qquad\qquad
\begin{tikzpicture}[scale=0.85]
  \node[font=\large] at (0,0) {$\otimes$};
  \foreach \r in {0.55,1.05,1.55}
    \draw[omThm,-Stealth] (\r,0) arc (360:10:\r);
  \node[omThm, font=\small] at (45:1.9) {$\vect B$};
  \node[font=\small] at (0,-2.1) {电流背离读者};
\end{tikzpicture}

{\small 直线导线周围的圆形磁感线，从正面看去
（$\odot$：电流垂直纸面向外；$\otimes$：电流垂直纸面向里）。用
右手握住导线，拇指沿电流方向，四指弯曲的方向即
$\vect B$ 的环绕方向。}
\end{omfigure}

\begin{example}[导线是一块弱磁铁]\label{ex:g11:magnetic-fields:wire}
距导线 $d = \qty{1.0}{cm}$ 处，导线中电流 $I = \qty{15}{A}$：
$B = \num{2e-7} \times 15 / 0.010 = \qty{0.30}{mT}$——相当于
六个地磁场，而这一电流足以使家用断路器跳闸。强磁场
需要更好的几何构型。
\end{example}

\begin{definition}[线圈与螺线管]\label{def:g11:magnetic-fields:solenoid}
绕制导线可使磁场集中。一个\emph{扁平
线圈}\index{线圈}（$N$ 匝绕于圆盘上）的行为类似一块薄磁铁，
一面为北极，另一面为南极；一个\emph{螺线管}\index{螺线管}——
即长圆柱形绕组，每米 $n$ 匝——其行为像条形磁铁。
\end{definition}

\begin{proposition}[螺线管内部的磁场]\label{prop:g11:magnetic-fields:solenoidfield}
在长螺线管内部、远离两端处，磁场是
\emph{均匀}：与轴线平行，且内部各点的大小相同，
其大小为 $B = \mu_0\, n\, I$，与
管径无关。管外磁场较弱。
\end{proposition}
\admitted

\begin{remark}
两式均在本科一年级分册中有严格推导，依据的是
磁场环流与所包围电流之间的关系。注意
真正起作用的是 $n$，即\emph{每米}匝数——密绕，而非加长。
\end{remark}

\begin{method}[右手定则两用]\label{met:g11:magnetic-fields:righthand}
\begin{itemize}
  \item \emph{导线}：拇指为电流方向；四指弯曲方向为磁感线的环绕方向。
  \item \emph{线圈、螺线管}：四指弯曲方向为线匝中电流的方向；
        拇指为内部 $\vect B$ 方向，指向北极面。
\end{itemize}
\end{method}

\begin{omfigure}
\begin{tikzpicture}[scale=0.9]
  \foreach \x in {0,0.65,...,3.9} {
    \node[font=\small] at (\x,0.85) {$\odot$};
    \node[font=\small] at (\x,-0.85) {$\otimes$};}
  \draw[-Stealth, omThm, thick] (-0.4,0.4) -- (4.3,0.4);
  \draw[-Stealth, omThm, thick] (-0.4,0) -- (4.3,0)
    node[right, font=\small] {$\vect B$};
  \draw[-Stealth, omThm, thick] (-0.4,-0.4) -- (4.3,-0.4);
  \draw[omExo, -Stealth, dashed]
    (4.5,0.7) .. controls (5.9,2.3) and (-2.0,2.3) .. (-0.6,0.7);
\end{tikzpicture}

{\small 螺线管纵剖图（$\odot$/$\otimes$：线匝
垂直穿过纸面）：内部沿轴线方向为均匀的 $B = \mu_0 n I$；
外部为较弱、分散的回流磁场（虚线）。}
\end{omfigure}

\begin{example}[螺线管的数值]\label{ex:g11:magnetic-fields:solenoid}
每厘米 $10$ 匝（$n = \qty{1000}{m^{-1}}$），且
$I = \qty{5.0}{A}$：$B = \num{1.26e-6} \times 1000 \times 5.0 \approx
\qty{6.3}{mT}$——相当于一百个地磁场，可用旋钮调节。
\end{example}

\section{电磁铁}

\begin{definition}[电磁铁]\label{def:g11:magnetic-fields:electromagnet}
一个\emph{电磁铁}\index{电磁铁}是绕在软铁芯上的螺线管：
铁芯在线圈磁场中被磁化，并增强该磁场，
通常增强 $100$ 至 $1000$ 倍。与永磁体不同，它
随电流切断而失去磁性。
\end{definition}

\begin{remark}[三种装置]\label{rem:g11:magnetic-fields:machines}
废品场的起重机用电磁铁吊起汽车，并且——
这正是关键——通过断开开关使其\emph{松开}。继电器的小
电磁铁吸合触点：以弱电流控制
强电流（\cref{ch:g11:circuits-and-power}）。核磁共振扫描仪
用超导螺线管在患者周围产生均匀的 \qty{1.5}{T} 磁场。
\end{remark}

各类磁场的数量级，从星际空间到中子星：
\begin{center}
\small
\begin{tabular}{lr}
星际空间 & $\sim \qty{e-10}{T}$ \\
地球表面的磁场 & \qty{5e-5}{T} \\
冰箱磁贴 & $\sim \qty{5e-3}{T}$ \\
钕磁铁，在其磁极处 & $\sim \qty{0.5}{T}$ \\
核磁共振螺线管 & \qty{1.5}{T} \\
中子星表面 & $\sim \qty{e8}{T}$ \\
\end{tabular}
\end{center}

\section{拉普拉斯力}

\begin{proposition}[拉普拉斯力]\label{prop:g11:magnetic-fields:laplace}
长为 $L$、通有电流 $I$ 的直导线，置于均匀
磁场 $\vect B$ 中，并与导线\emph{垂直}，受到
\emph{拉普拉斯力}\index{拉普拉斯力}，大小为 $F = I\,L\,B$，
方向与导线及 $\vect B$ 均垂直：右手展平，拇指沿电流方向，
伸直的四指沿 $\vect B$ 方向——掌心推动的方向即
$\vect F$ 的方向。与磁场\emph{平行}的导线不受力。
\end{proposition}
\admitted

\begin{remark}[跳动的导轨]\label{rem:g11:magnetic-fields:rail}
本科一年级分册由磁场对每个运动
电荷的作用力推出此结论。实验中，一根铜棒横放在两根水平
导轨上，位于磁铁两极之间（磁场垂直纸面向里，
见图）；闭合开关，铜棒便沿导轨飞速射出。
调换电流方向或翻转磁铁，铜棒便向另一方向射出。
\end{remark}

\begin{omfigure}
\begin{tikzpicture}[scale=0.95]
  \foreach \x in {0.6,1.5,2.9,3.8}
    \foreach \y in {0.35,1.05}
      \node[omThm, font=\small] at (\x,\y) {$\otimes$};
  \node[omThm, font=\small] at (4.35,1.05) {$\vect B$};
  \draw[very thick] (0,0) -- (4.6,0);
  \draw[very thick] (0,1.4) -- (4.6,1.4);
  \draw[omDef, line width=2.5pt] (2.1,-0.12) -- (2.1,1.52);
  \draw[-Stealth, omDef, thick] (0.1,1.6) -- (1.6,1.6)
    node[midway, above, font=\small] {$I$};
  \draw[-Stealth, omDef, thick] (2.32,1.15) -- (2.32,0.45);
  \draw[-Stealth, omDef, thick] (1.6,-0.22) -- (0.1,-0.22);
  \draw[-Stealth, omMeth, very thick] (2.1,0.7) -- (3.4,0.7)
    node[pos=0.85, above, font=\small] {$\vect F$};
\end{tikzpicture}

{\small 导轨实验：电流沿铜棒向下，磁场垂直纸面向里
（$\otimes$）——拉普拉斯力 $F = ILB$ 驱动铜棒
沿导轨运动。}
\end{omfigure}

\begin{example}[拉普拉斯力的数值]\label{ex:g11:magnetic-fields:laplace}
一根长为 $L = \qty{5.0}{cm}$ 的导体棒，通有 $I = \qty{8.0}{A}$ 的电流，
垂直置于 $B = \qty{0.50}{T}$ 的磁场中：$F = 8.0 \times 0.050 \times 0.50 = \qty{0.20}{N}$
——相当于质量为 \qty{20}{g} 的物体所受的重力，而这段导线本身只重几克。
\end{example}

\begin{remark}[直流电动机]\label{rem:g11:magnetic-fields:motor}
将导线在磁场中弯成线圈：与
$\vect B$ 垂直的两边电流方向相反，所受拉普拉斯力也相反，
——构成一对使线圈转动的力偶。转过半圈后，这对力会把线圈反向拨回，
因此需要一个旋转开关（\emph{换向器}\index{换向器}）每转过半圈便
使电流反向。每一台风扇和电钻，都是这种线圈的多匝版本。
\end{remark}

\section{练习}

\begin{exercise}[$\star$]\label{exo:g11:magnetic-fields:1}
一根条形磁铁被从中间锯成两段。每一段各带
什么磁极？重新拼合后，新切开的断面会相互吸引
还是相互排斥？能否通过任何切割方式分离出北极？
\end{exercise}

\begin{exercise}[$\star$]\label{exo:g11:magnetic-fields:2}
判断下列说法是否正确，并各说明理由：(a) 两条磁感线
在磁场强的地方相交；(b) 磁铁外部的磁感线由
北极指向南极；(c) 每一条磁感线都是
闭合曲线；(d) 小磁针静止时与过它的磁感线垂直。
\end{exercise}

\begin{exercise}[$\star$]\label{exo:g11:magnetic-fields:3}
求距直导线 \qty{2.0}{cm} 处的磁场，该导线通有
\qty{10}{A} 的电流。该磁场是地磁场 \qty{50}{\micro T} 的多少倍？
\end{exercise}

\begin{exercise}[$\star$]\label{exo:g11:magnetic-fields:4}
一螺线管由 $800$ 匝导线均匀绕在 \qty{40}{cm} 长度上，
通有 $I = \qty{1.5}{A}$ 的电流。求 $n$，再求内部的 $B$；若电流加倍，
$B$ 变为多少？
\end{exercise}

\begin{exercise}[$\star$]\label{exo:g11:magnetic-fields:5}
一段长为 \qty{10}{cm} 的导线，通有 \qty{5.0}{A} 的电流，
垂直置于 \qty{0.20}{T} 的磁场中。求拉普拉斯力；
该力与质量为多大的物体所受的重力相等？
\end{exercise}

\begin{exercise}[$\star\star$]\label{exo:g11:magnetic-fields:6}
徒步者所在处的磁偏角为西偏 $\ang{8}$。她按
“指南针指北”的方向走 \qty{2.0}{km}，求她偏离真北多远、偏向哪一侧，
以及若要真正走向真北，应取怎样的方位？
\end{exercise}

\begin{exercise}[$\star\star$]\label{exo:g11:magnetic-fields:7}
距一根通有 \qty{20}{A} 电流的直导线多远处，其磁场
恰好等于地磁场的 \qty{50}{\micro T}？这对靠近通电电缆时
依赖指南针意味着什么？
\end{exercise}

\begin{exercise}[$\star\star$]\label{exo:g11:magnetic-fields:8}
设计一螺线管，使之产生 $B = \qty{10}{mT}$，电流为 $I = \qty{4.0}{A}$：
求所需的 $n$，再求长为 \qty{25}{cm} 的线圈所需的匝数。
\end{exercise}

\begin{exercise}[$\star\star$]\label{exo:g11:magnetic-fields:9}
在导轨实验中，铜棒（质量 \qty{20}{g}，两导轨之间的长度
$L = \qty{12}{cm}$）通有 \qty{10}{A} 的电流，置于
\qty{0.50}{T} 的垂直磁场中。求拉普拉斯力，再求
铜棒的初始加速度，并与 $g$ 比较。
\end{exercise}

\begin{exercise}[$\star\star$]\label{exo:g11:magnetic-fields:10}
一根水平导线，质量为 \qty{5.0}{g}，长为 \qty{20}{cm}，
置于与之垂直的水平磁场 $B = \qty{0.10}{T}$ 中。要使
拉普拉斯力与导线重力平衡、导线悬浮，需多大的电流？
该力应指向哪一方向？
\end{exercise}

\begin{exercise}[$\star\star$]\label{exo:g11:magnetic-fields:11}
给出拉普拉斯力的方向（若为零，请说明）：
(a) 电流沿纸面向右，$\vect B$ 垂直纸面向里；
(b) 电流沿纸面向上，$\vect B$ 垂直纸面向外；
(c) 电流与 $\vect B$ 平行。可用平展的右手判断。
\end{exercise}

\begin{exercise}[$\star\star\star$]\label{exo:g11:magnetic-fields:12}
一根水平导线沿地理南–北方向铺设，位于一只指南针上方 \qty{3.0}{cm} 处；
该处地磁场的水平分量为 \qty{20}{\micro T}。
导线中通有 $I = \qty{5.0}{A}$ 的电流：求其在指南针处产生的磁场，
给出其方向，再求小磁针的偏转角（磁场按矢量相加）。
\end{exercise}

\begin{exercise}[$\star\star\star$]\label{exo:g11:magnetic-fields:13}
废料场电磁铁是一个线圈，长 \qty{20}{cm} 上绕有 $400$ 匝，
电阻为 \qty{2.0}{\ohm}，接在 \qty{12}{V} 的电源上；其铁芯
使裸线圈的磁场放大为 $100$ 倍。求电流、裸线圈的磁场，
加铁芯后的磁场，以及线圈耗散的功率
（见\cref{ch:g11:circuits-and-power}）。为什么废料在开关断开的
瞬间便会落下？
\end{exercise}

\begin{exercise}[$\star\star\star$]\label{exo:g11:magnetic-fields:14}
电动机的矩形线框有 $N = 50$ 匝，通有
$I = \qty{2.0}{A}$ 的电流；其长为 \qty{5.0}{cm} 的两边
与磁场 $B = \qty{0.30}{T}$ 垂直。求这两边各自所受的力，
为何两力方向相反？它们对线框有何作用？
另外两边为何有时不受力？为何必须
每转过半圈就将电流反向？
\end{exercise}

\begin{exercise}[$\star\star\star$]\label{exo:g11:magnetic-fields:15}
距离通有 \qty{100}{A} 电流的电缆多近，其磁场才能达到
\qty{1}{T}？将其与电缆的毫米级半径相比较，
并归纳：特斯拉量级的磁场为何要用
螺线管而非单根导线来产生。再求：中子星表面磁场为
\qty{e8}{T}，比地球的磁场高多少个数量级？
\end{exercise}

\section{问题：指南针、起重机与磁共振成像}

\begin{problem}[{周末问题——三种磁体各显身手：指南针横渡大洋、
电磁铁吊起汽车、磁共振成像的螺线管
——同一磁场跨越十个数量级，承担三项任务}]
\label{pb:g11:magnetic-fields:1}
同一矢量 $\vect B$ 可用 \qty{50}{\micro T} 引导帆船，
用 \qty{1}{T} 吊起废铁，用 \qty{1.5}{T} 为膝关节成像——
——三台机器分别对应本章的三个公式。

\medskip
\noindent\textbf{第一部分——指南针。}
船所在位置处，地磁场的水平分量为
\qty{20}{\micro T}，竖直分量为 \qty{44}{\micro T}；
磁偏角为西偏 $\ang{6}$。
\begin{enumerate}
  \item 为什么指南针的指针究竟会指向（磁）北方？
  \item 求总磁场的大小及其与水平方向的夹角。
  \item 若按“指南针北向”航行 \qty{30}{km}，
        船最终将位于真北航线以西多远？
  \item 钢铁货物使指针再向西偏 $\ang{4}$，问题同上。
  \item 靠近地磁极时指南针失效，哪一个分量应负其责，
        为什么？
\end{enumerate}

\medskip
\noindent\textbf{第二部分——起重机。}
废料场电磁铁是一个螺线管，共 $600$ 匝，
密绕于 \qty{30}{cm} 长度上，电阻为 \qty{1.5}{\ohm}，由 \qty{12}{V}
电源供电；其铁芯使裸线圈的磁场放大为 $50$ 倍。
\begin{enumerate}[resume]
  \item 求匝数密度 $n$。
  \item 求线圈中的电流（见\cref{ch:g11:circuits-and-power}）。
  \item 求裸线圈的磁场 $\mu_0 n I$。
  \item 求加铁芯后的磁场，并与钕磁铁的 \qty{0.5}{T} 比较。
  \item 求线圈耗散的功率，再求
        一个 \qty{10}{\minute} 吊运班次所需的能量。
  \item 同样强度的永磁体确实存在，那为何仍用电磁铁？
\end{enumerate}

\medskip
\noindent\textbf{第三部分——绞车电动机。}
起重机的绞车是一台直流电动机：矩形线框共 $N = 100$ 匝，
长为 \qty{4.0}{cm} 的边与 $B = \qty{0.25}{T}$ 垂直，
通有 $I = \qty{3.0}{A}$ 的电流。
\begin{enumerate}[resume]
  \item 求每一条与 $\vect B$ 垂直的边所受的拉普拉斯力。
  \item 两力大小相等、方向相反，但其效果并不会
        相互抵消。它们产生什么效果？为何必须\emph{相反}？
  \item 当另外两边与 $\vect B$ 平行时，它们受到什么力？
  \item 转过半圈后，同样的这对力将抵消转动。
        换向器起何作用，又在何时动作？
  \item 给出三种设计改动，各自都能使这对力加倍。
\end{enumerate}

\medskip
\noindent\textbf{第四部分——磁共振成像。}
扫描仪的螺线管以电流
$I = \qty{500}{A}$ 维持 \qty{1.5}{T} 的均匀磁场。
\begin{enumerate}[resume]
  \item \qty{1.5}{T} 相当于地球磁场的多少倍？
  \item 求\emph{不加}任何铁芯时所需的匝数密度 $n$。
  \item 若绕组的电阻为 \qty{0.20}{\ohm}，
        则其耗散功率为多少？实际绕组凭借何种性质避免了这一损耗，
        代价是极低的温度？
  \item 总结：把船、起重机和磁共振成像的磁场
        标在数量级阶梯上，从星际空间
        （\qty{e-10}{T}）到中子星（\qty{e8}{T}）；
        各台机器之间变化的是物理规律——还是电流大小？
\end{enumerate}
\end{problem}
